\subsection{分解群与惯性群}

定义 2. 设 A' 是环，G 是作用于 A' 的群（§ 1，第 9 节）。给定 A' 的素理想 p'，由满足 σ ∈ G 且 σ.p' = p' 的元素 σ 组成的子群称为 p' （关于 G）的分解群，记作 \(G^{\mathrm{Z}}(p')\)。A' 中在 \(G^{\mathrm{Z}}(p')\) 下不变的元素组成的环称为 p' （关于 G）的分解环，记作 \(A^{\mathrm{Z}}(p')\) (*)。

  当不致歧义时，我们常用 \(\mathcal{G}^{\mathbf{Z}}\) 与 \(\mathbf{A}^{\mathbf{Z}}\) 分别代替 \(\mathcal{G}^{\mathbf{Z}}(\mathfrak{p}')\) 与 \(\mathbf{A}^{\mathbf{Z}}(\mathfrak{p}')\) 。

对所有 \(\sigma \in \mathcal{G}^{\mathbb{Z}}(\mathfrak{p}')\) ，我们也以 \(z \mapsto \sigma.z\) 表示环 \(\mathbf{A}' / \mathfrak{p}'\) 的自同态，它是由自同态 \(x \mapsto \sigma.x\) （\(\mathbf{A}'\) 的）取商得到的；显然群 \(\mathcal{G}^{\mathbb{Z}}(\mathfrak{p}')\) 以这种方式作用于环 \(\mathbf{A}' / \mathfrak{p}'\) 。

定义 3. 用定义 2 的记号，\(\mathcal{G}^{\mathrm{Z}}(\mathfrak{p}')\) 中由那些 \(\sigma\) （它们使自同态 \(z \mapsto \sigma.z\) ——\(\mathrm{A}' / \mathfrak{p}'\) 的——为恒同）组成的子群，称为 \(\mathfrak{p}'\) （关于 \(\mathcal{G}\) ）的惯性群，记作 \(\mathcal{G}^{\mathrm{T}}(\mathfrak{p}')\) （或 \(\mathcal{G}^{\mathrm{T}}\) ）。\(\mathrm{A}'\) 中在 \(\mathcal{G}^{\mathrm{T}}(\mathfrak{p}')\) 下不变的元素组成的环称为 \(\mathfrak{p}'\) （关于 \(\mathcal{G}\) ）的惯性环，记作 \(\mathrm{A}^{\mathrm{T}}(\mathfrak{p}')\) （或 \(\mathrm{A}^{\mathrm{T}}\) ）（ \(\dagger\) ）。

若 A 是由 G 的不变元组成的 A' 的子环，则显然

\[
\mathbf {A} \subset \mathbf {A} ^ {\mathbf {Z}} (\mathfrak {p} ^ {\prime}) \subset \mathbf {A} ^ {\mathbf {T}} (\mathfrak {p} ^ {\prime}) \subset \mathbf {A} ^ {\prime}.\tag{1}
\]

由定义 2 与定义 3 可得，对所有 \(\rho \in \mathcal{G}\) ，

\[
\mathcal {G} ^ {\mathrm{Z}} (\rho \cdot p ^ {\prime}) = \rho \mathcal {G} ^ {\mathrm{Z}} (p ^ {\prime}) \rho^ {- 1}, \quad \mathcal {G} ^ {\mathrm{T}} (\rho \cdot p ^ {\prime}) = \rho \mathcal {G} ^ {\mathrm{T}} (p ^ {\prime}) \rho^ {- 1}.\tag{2}
\]

若对所有 \(\sigma \in \mathcal{G}^{\mathbb{Z}}(\mathfrak{p}')\) ，\(\bar{\sigma}\) 表示自同构 \(z \mapsto \sigma. z\) （\(\mathrm{A}' / \mathfrak{p}'\) 的），则 \(\sigma \mapsto \bar{\sigma}\) 是从 \(\mathcal{G}^{\mathbb{Z}}\) 到群 \(\Gamma_0\) 的同态（称为典范同态），该群由

\(A^{\prime}/p^{\prime}\) 的使 \(A^{Z}/(p^{\prime} \cap A^{Z})\) （典范地等同于 \(A^{\prime}/p^{\prime}\) 的子环）的元素不变的诸自同构组成；按定义 \(\mathcal{G}^{\mathrm{T}}(p^{\prime})\) 是这个典范同态的核；因而 \(G^{T}\) 是 \(G^{Z}\) 的正规子群。若 \(k^{\prime}\) 是 \(A^{\prime}/p^{\prime}\) 的分式域，则 \(A^{\prime}/p^{\prime}\) 的每个自同构可唯一地扩张为 \(k^{\prime}\) 的自同构，于是 \(\sigma \mapsto \bar{\sigma}\) 可看作从 \(\mathcal{G}^{Z}(p^{\prime})\) 到 \(k^{\prime}\) 的自同构群的同态。最后注意，由于 \(G^{T}\) 在 \(G^{Z}\) 中正规，\(A^{T}\) 在 \(G^{Z}\) 下稳定。

引理 3. 设 \(\mathbf{A}'\) 是环，\(\mathcal{G}\) 是作用于 \(\mathbf{A}'\) 的群，\(\mathbf{A}\) 是 \(\mathcal{G}\) 的不变元组成的环，\(\mathfrak{p}'\) 是 \(\mathbf{A}'\) 的素理想，\(\mathbf{S}\) 是 \(\mathbf{A}\) 的与 \(\mathfrak{p}'\) 不相交的乘性子集。则 \(\mathcal{G}^{\mathrm{Z}}(\mathrm{S}^{-1}\mathfrak{p}') = \mathcal{G}^{\mathrm{Z}}(\mathfrak{p}')\) ，\(\mathcal{G}^{\mathrm{T}}(\mathrm{S}^{-1}\mathfrak{p}') = \mathcal{G}^{\mathrm{T}}(\mathfrak{p}')\) ，且若 \(\mathcal{G}\) 局部有限，则 \(\mathrm{S}^{-1}\mathrm{A}^{\mathrm{Z}}(\mathfrak{p}') = \mathrm{A}^{\mathrm{Z}}(\mathrm{S}^{-1}\mathfrak{p}')\) 且 \(\mathrm{S}^{-1}\mathrm{A}^{\mathrm{T}}(\mathfrak{p}') = \mathrm{A}^{\mathrm{T}}(\mathrm{S}^{-1}\mathfrak{p}')\) 。

由于 S 的元素在 \(\mathcal{G}\) 下不变，显然，若 \(\sigma. p' = p'\) ，则也有 \(\sigma. (S^{-1}p') = S^{-1}p'\) 。反之，设 \(\sigma \in \mathcal{G}\) 满足 \(\sigma. (S^{-1}p') = S^{-1}p'\) ；则若 \(x \in p'\) ，有 \((\sigma.x)/1 \in S^{-1}p'\) ，因此存在 \(s \in S\) 使得 \(s(\sigma.x) \in p'\) ，从而 \(\sigma.x \in p'\) ，因为 \(p'\) 素且 \(s \notin p'\) ；这就证明了 \(\sigma. p' \subset p'\) ，类似地可证 \(\sigma^{-1}. p' \subset p'\) ，因而 \(\sigma. p' = p'\) ，即 \(\sigma \in \mathcal{G}^{\mathrm{Z}}(p')\) 。若 \(\sigma \in \mathcal{G}^{\mathrm{T}}(p')\) ，则 \(\sigma.x - x \in p'\) 对所有 \(x \in A'\) 成立，因而对所有 \(s \in S\) 也有

\[
\sigma . (x / s) - (x / s) = (\sigma . x - x) / s \in \mathbf {S} ^ {- 1} \mathfrak {p} ^ {\prime}
\]

从而 \(\sigma \in \mathcal{G}^{\mathrm{T}}(\mathrm{S}^{-1}\mathfrak{p}')\) 。反之，设 \(\sigma \in \mathcal{G}^{\mathrm{T}}(\mathrm{S}^{-1}\mathfrak{p}')\) ；则对所有 \(x \in \mathbf{A}'\) ，\(\sigma.(x/1) - (x/1) \in \mathrm{S}^{-1}\mathfrak{p}'\) ，因此存在 \(s \in \mathbf{S}\) 使得 \(s(\sigma.x - x) \in \mathfrak{p}'\) ，于是与上面一样 \(\sigma.x - x \in \mathfrak{p}'\) ，这就证明了 \(\sigma \in \mathcal{G}^{\mathrm{T}}(\mathfrak{p}')\) 。最后两个断言由 § 1，第 9 节，命题 23 得出。

定理 2. 设 A' 是环，G 是作用于 A' 的有限群，A 是 G 的不变元组成的环，于是 A' 在 A 上是整的（§ 1，第 9 节，命题 22）。

\begin{enumerate}
\def\labelenumi{(\roman{enumi})}
\item
  给定两个素理想 \(\mathfrak{p}'\) 、\(\mathfrak{q}'\) （\(\mathbf{A}'\) 的，位于同一素理想 \(\mathfrak{p}\) （\(\mathbf{A}\) 的）之上），存在 \(\sigma \in \mathcal{G}\) 使得 \(\mathfrak{q}' = \sigma \cdot \mathfrak{p}'\) ；换言之，\(\mathcal{G}\) 传递地作用于 \(\mathbf{A}'\) 的位于 \(\mathfrak{p}\) 之上的素理想组成的集。
\item
  设 \(\mathfrak{p}'\) 是 \(\mathbf{A}', \mathfrak{p} = \mathfrak{p}' \cap \mathbf{A}\) 的素理想，并设 \(k\) （相应地 \(k'\) ）是 \(\mathbf{A}/\mathfrak{p}\) （相应地 \(\mathbf{A}'/\mathfrak{p}'\) ）的分式域。则 \(k'\) 是 \(k\) 的拟 Galois 扩张 (*)，且典范同态 \(\sigma \mapsto \bar{\sigma}\) （从 \(\mathcal{G}^{\mathrm{Z}}(\mathfrak{p}')\) 到群 \(\Gamma\) ——由 \(k\) -自同构（\(k'\) 的）组成的群——的）取商后定义了 \(\mathcal{G}^{\mathrm{Z}}(\mathfrak{p}') / \mathcal{G}^{\mathrm{T}}(\mathfrak{p}')\) 到 \(\Gamma\) 上的同构。
\setcounter{enumi}{0}
\item
  若 \(x \in \mathfrak{q}'\) ，则 \(\prod_{\sigma \in \mathcal{G}} \sigma. x \in \mathfrak{q}' \cap A = \mathfrak{p} \subset \mathfrak{p}'\) ；于是存在 \(\sigma \in \mathcal{G}\) 使得 \(\sigma. x \in \mathfrak{p}'\) ，即 \(x \in \sigma^{-1}. \mathfrak{q}'\) 。由此可断定 \(\mathfrak{q}' \subset \bigcup_{\sigma \in \mathcal{G}} \sigma. \mathfrak{p}'\) ，因而（由于 \(\mathcal{G}\) 有限且诸 \(\sigma. \mathfrak{p}'\) 为素理想）存在 \(\sigma \in \mathcal{G}\) 使得 \(\mathfrak{q}' \subset \sigma. \mathfrak{p}'\) （第 II 章，§ 1，第 1 节，命题 2）；由于 \(\mathfrak{q}'\) 与 \(\sigma. \mathfrak{p}'\) 都位于 \(\mathfrak{p}, \mathfrak{q}' = \sigma. \mathfrak{p}'\) （第 1 节，命题 1 的推论 1）。
\end{enumerate}

\iffalse
\[
k ^ {\prime}
\]

\[
k,
\]

\[
\bar {x} \in A ^ {\prime} / p ^ {\prime}
\]

\[
k [ \mathbf {X} ]
\]

\[
\mathbf {A} ^ {\prime} / \mathfrak {p} ^ {\prime}
\]

\[
\boldsymbol {x} \in \mathbf {A} ^ {\prime}
\]

\[
\mathrm{V}, \S 6,
\]

\[
\overline {{x}};
\]
\fi

(ii) 为证 \(k'\) 是 \(k\) 的拟 Galois 扩张，只需证明每个元素 \(\bar{x}\) （\(\mathbf{A}'/\mathfrak{p}'\) 的）都是某个多项式 \(P\) （\(k[\mathbf{X}]\) 中的）的根，且该多项式的全部根都在 \(\mathbf{A}'/\mathfrak{p}'\) 中（《代数》，第 V 章，§ 6，第 3 节，命题 9 的推论 3）。现设 \(x \in \mathbf{A}'\) 是类 \(\bar{x}\) 的一个代表；多项式

\(Q(X) = \prod_{\sigma \in \mathcal{G}} (X - \sigma.x)\) 的全部系数都在 A 中；设 \(P(X)\) 是 \((A/p)[X]\) 中的多项式，其系数是 Q 的诸系数在典范同态 \(\pi: A \to A/p\) 下的像。由于 \(\pi\) 可看作典范同态 \(\pi': A' \to A'/p'\) 在 A 上的限制，可见 P 是 \((A'/p')[X]\) 中诸线性因子 \(X - \pi'(\sigma.x)\) 之积，因而解决了问题，因为 \(\bar{x} = \pi'(x)\) 。

显然，对所有 \(\sigma \in \mathcal{G}^{\mathbb{Z}}\) ，\(\bar{\sigma}\) 是 \(k\) -自同构（\(k'\) 的）；剩下要验证 \(\sigma \mapsto \bar{\sigma}\) 把 \(\mathcal{G}^{\mathbb{Z}}\) 映到全部 \(k\) -自同构（\(k'\) 的）组成的群上。记 \(\mathrm{S} = \mathrm{A} - \mathfrak{p}\) ；\(k\) 与 \(k'\) 不因把 \(\mathrm{A}'\) 与 \(\mathfrak{p}'\) 分别换成 \(\mathrm{S}^{-1}\mathrm{A}'\) 与 \(\mathrm{S}^{-1}\mathfrak{p}'\) 而改变（依据 § 1，第 9 节，命题 23 及关系式 \(\mathrm{S}^{-1}\mathfrak{p}' \cap \mathrm{S}^{-1}\mathrm{A} = \mathrm{S}^{-1}(\mathrm{A} \cap \mathfrak{p}') = \mathrm{S}^{-1}\mathfrak{p}\) ，第 II 章，§ 2，第 4 节）；由引理 3，\(\mathcal{G}^{\mathbb{Z}}\) 及其在 \(k'\) 上的作用也不改变；因此只需考察 \(\mathfrak{p}\) 极大的情形，此时我们知道 \(\mathfrak{p}'\) 也是极大的（第 1 节，命题 1），因而 \(k'\) 的每个元素都形如 \(\pi'(x)\) ，其中 \(x \in \mathrm{A}'\) ；上面已经看到，这样的元素是 \(k[\mathrm{X}]\) 中某个次数 \(\leqslant \operatorname{Card}(\mathcal{G})\) 的多项式的根。由于 \(k\) 的每个有限可分扩张都有本原元素（《代数》，第 V 章，§ 7，第 7 节，命题 12 及 § 11，第 4 节，命题 4），可见 \(k\) 的每个含于 \(k'\) 的有限可分扩张的次数都 \(\leqslant \operatorname{Card}(\mathcal{G})\) ，由此推出 \(k_s'\) ——\(k\) 的、含于 \(k'\) 的最大可分扩张——（《代数》，第 V 章，§ 7，第 6 节，命题 11）的次数 \(\leqslant \operatorname{Card}(\mathcal{G})\) （《代数》，第 V 章，§ 3，第 2 节，注 2）。设 \(y \in \mathrm{A}'\) 是使 \(\pi'(y)\) 为 \(\mathbf{k}_s'\) 的本原元素的元素。按定义，理想 \(\sigma .\mathfrak{p}'\) （对 \(\sigma \in \mathcal{G} - \mathcal{G}^{\mathbb{Z}}\) ）都是极大的且异于 \(\mathfrak{p}'\) ；因此存在 \(x \in \mathrm{A}'\) 使得 \(x \equiv y (\text{mod. } \mathfrak{p}')\) 且 \(x \in \sigma^{-1}\mathfrak{p}'\) （对 \(\sigma \in \mathcal{G} - \mathcal{G}^{\mathbb{Z}}\) ）（第 II 章，§ 1，第 2 节，命题 5）。在这样的约定下，设 \(u\) 是一个 \(k\) -自同构（\(k'\) 的），并设 \(\mathrm{P}(\mathrm{X}) = \prod_{\sigma \in \mathcal{G}} (\mathrm{X} - \pi'(\sigma .x))\) ；由于 \(\pi'(x)\) 是 P 的根且 \(\mathrm{P} \in k[\mathrm{X}]\) ，\(u(\pi'(x))\) 也是 P 在 \(k'\) 中的根，因而存在 \(\tau \in \mathcal{G}\) 使得

\[
u (\pi^ {\prime} (x)) = \pi^ {\prime} (\tau . x);
\]

  但 \(u(\pi'(x)) \neq 0\) ，而对 \(\sigma \in \mathcal{G} - \mathcal{G}^{\mathbf{Z}}\) ，\(\sigma.x \in \mathfrak{p}'\) ，从而 \(\pi'(\sigma.x) = 0\) ；由此断定必有 \(\tau \in \mathcal{G}^{\mathbf{Z}}\) 。而由于 \(u\) 与 \(\bar{\tau}\) 在本原元素 \(\pi'(y) = \pi'(x)\) （\(k_s'\) 的）上取相同的值，它们在 \(k_s'\) 上一致；又由于 \(k'\) 是 \(k_s'\) 的纯不可分扩张，它们在 \(k'\) 上一致。

推论. 在定理 2 的假设与记号下，设 \(f_1, f_2\) 是从 \(\mathbf{A}'\) 到域 \(\mathbf{L}\) 的两个同态，且它们在 \(\mathbf{A}'\) 上的限制相同。则存在 \(\sigma \in \mathcal{G}\) 使得

\[
f _ {2} (x) = f _ {1} (\sigma . x)
\]

对所有 \(x \in A'\) 。

设 \(p_{i}'\) 是 \(f_{i}\) ( \(i = 1, 2\) )的核，它是 \(A'\) 的素理想；按假设 \(p_{1}' \cap A = p_{2}' \cap A\) ，且这个交是素理想 \(p\) （\(A\) 的）；因此存在 \(\tau \in \mathcal{G}\) 使得 \(\tau. p_{2}' = p_{1}'\) （定理 2 (i)）；把 \(f_{1}\) 换成同态 \(x \mapsto f_{1}(\tau.x)\) ，便可设 \(p_{2}' = p_{1}'\) （把这个理想记作 \(p'\) ）。取商后，由 \(f_{1}\) 与 \(f_{2}\) 得到两个单同态 \(f_{1}', f_{2}'\) （从 \(A'/p'\) 到 \(L\) 的），它们因而分别扩张为两个单同态 \(f_{1}''\) 、\(f_{2}''\) （从 \(k'\) ——\(A'/p'\) 的分式域——到 \(L\) 的）。由于 \(k'\) 是 \(k\) 的拟 Galois 扩张，\(k_{1}'' = f_{1}''(k')\) 与 \(k_{2}'' = f_{2}''(k')\) 也是（ \(k\) 等同于 \(L\) 的子域），且由于存在 \(k\) -同构（从 \(k_{1}''\) 到 \(k_{2}''\) 上的），故 \(k_{1}'' = k_{2}''\) （《代数》，第 V 章，§ 6，命题 6）。于是 \(f_{1}'' \circ f_{2}''\) 是 \(k\) -自同构（\(k'\) 的）；由定理 2 (ii)，它因而形如 \(\bar{\sigma}\) ，其中 \(\sigma \in \mathcal{G}^{Z}(p')\) 。特别地，对所有 \(x \in A'\) ，元素 \(f_{2}(x)\) 与 \(f_{1}(\sigma.x)\) 相等。

\textbf{注}\par

\begin{enumerate}
\def\labelenumi{(\arabic{enumi})}
\item
  注意，在定理 2 的假设下，若 \(k'\) 在 \(k\) 上不可分，则 \(k'\) 在 \(k\) 上可以是无限的（习题 9）。
\item
  显然 \(k'\) 是 \(k\) 的 Galois 扩张，若域 \(k\) 是完满的。此时它在 \(k\) 上是有限的。
\end{enumerate}

命题 4. 设 A' 是环，G 是作用于 A' 的有限群，H 是 G 的子群，A 与 B 分别是 G 与 H 的不变元组成的环，p' 是 A' 的素理想；记 p = A ∩ p' ，p(B) = B ∩ p' 。

\begin{enumerate}
\def\labelenumi{(\roman{enumi})}
\item
  为使 \(\mathcal{H}\) 包含在分解群 \(\mathcal{G}^{\mathrm{Z}}(\mathfrak{p}')\) 中，必须而且只需 \(\mathfrak{p}'\) 是 A 的位于 \(\mathfrak{p}(\mathbf{B})\) 之上的唯一素理想。
\item
  若 \(\mathcal{H}\) 包含 \(\mathcal{G}^{\mathrm{Z}}(\mathfrak{p}')\) ，则下列条件成立：
\end{enumerate}

\begin{enumerate}
\def\labelenumi{(\alph{enumi})}
\item
  环 A/p 与 B/p(B) 有相同的分式域。
\item
  局部环 \(\mathbf{B}_{\mathfrak{p}(\mathbf{B})}\) 的极大理想等于 \(\mathfrak{p}\mathbf{B}_{\mathfrak{p}(\mathbf{B})}\) 。
\end{enumerate}

\begin{enumerate}
\def\labelenumi{(\roman{enumi})}
\setcounter{enumi}{2}
\item
  进一步设 \(\mathbf{A}'\) 是整环且 \(\bigcap_{n\geqslant 0}\mathfrak{p}^n\mathbf{A}_{\mathfrak{p}'} = 0\) ；则 (ii) 的条件 (a) 与 (b) 蕴含 \(\mathcal{G}^{\mathrm{Z}}(\mathfrak{p}')\) 使 \(\mathbf{B}\) 的元素不变。

\setcounter{enumi}{0}
\item
  由定理 2 (i)，A' 的位于 p(B) 之上的素理想正是形如 \(\sigma. \mathfrak{p}'\) 的理想，其中 \(\sigma \in \mathcal{H}\) ；由此立得 (i)。

\item
  记 \(S = A - p\) ；我们知道 \(\mathcal{G}\) 与 \(\mathcal{H}\) 在 \(S^{-1}A'\) 中的不变元环分别是 \(S^{-1}A\) 与 \(S^{-1}B\) （§ 1，第 9 节，命题 23），且 \(\mathcal{G}^{\mathbb{Z}}(S^{-1}p') = \mathcal{G}^{\mathbb{Z}}(p')\) （引理 3）；最后有 \(S^{-1}p(B) = S^{-1}p' \cap S^{-1}B\) （第 II 章，§ 2，第 4 节），素理想 \(S^{-1}p(B)\) （环 \(S^{-1}B\) 的）的局部环典范同构于 \(B_{p(B)}\) ，其剩余域同构于 \(B/p(B)\) 的分式域（第 II 章，§ 2，第 5 节，命题 11）。因此，证明 (ii) 时可只考察 \(p\) 极大的情形。为证 (a)，只需证明
\end{enumerate}

\[
\mathbf {B} = \mathbf {A} + \mathfrak {p} (\mathbf {B})\tag{3}
\]

因为这将表明域 A/p 与 B/p(B) 典范同构。由定理 2，A' 的位于 p 之上的素理想只有有限个，而由第 1 节的定理 1，B 的每个素理想之上至少有 A' 的一个素理想；由此推出 B 的位于 p 之上的素理想只有有限个；设 \(n_{j}\) ( \(1 \leqslant j \leqslant r\) )表示这些理想中 \(\neq p(B)\) 的那些。设 x 是 B 的元素；由于理想 \(p(B)\) 与 \(n_{j}\) 都是极大的（第 1 节，命题 1），存在 \(y \in B\) 使得 \(y \equiv x \pmod{p(B)}\) 且 \(y \in n_{j}\) （对 \(1 \leqslant j \leqslant r\) ）（第 II 章，§ 1，第 2 节，命题 5）。设 \(y_{1} = y, y_{2}, \ldots, y_{q}\) 是 y 在 G 下轨道中的不同元素；显然

\[
z = y _ {1} + y _ {2} + \dots + y _ {q} \in \mathrm{A},
\]

而为证 (3)，只需证明 \(y_i \in \mathfrak{p}'\) （对 \(i \geqslant 2\)），因为由此将导出 \(z - y \in \mathfrak{p}' \cap B = \mathfrak{p}(B)\) ，从而 \(x \in A + \mathfrak{p}(B)\) ，因为 \(x \equiv y \pmod{\mathfrak{p}(B)}\) 。现设 \(i \geqslant 2\) 且 \(\sigma \in \mathcal{G}\) 满足 \(\sigma.y = y_i\) ；我们证明 \(\sigma^{-1}. \mathfrak{p}'\) 不位于 \(\mathfrak{p}(B)\) 之上。否则将存在 \(\tau \in \mathcal{H}\) 使得 \(\sigma^{-1}. \mathfrak{p}' = \tau. \mathfrak{p}'\) （定理 2 (i)），于是 \((\tau^{-1}\sigma^{-1}). \mathfrak{p}' = \mathfrak{p}'\) ，换言之 \(\tau^{-1}\sigma^{-1} \in \mathcal{G}^{Z} \subset \mathcal{H}\) （按假设），从而 \(\sigma \in \mathcal{H}\) ；但由于 \(y \in B\) 且 \(\sigma.y \neq y\) ，这是荒谬的。由此断定 \(\sigma^{-1}. \mathfrak{p}'\) 位于某个理想 \(n_j\) 之上，且由于按构造 \(y \in n_j\) ，当然有 \(y \in \sigma^{-1}. \mathfrak{p}'\) ，即 \(y_i = \sigma.y \in \mathfrak{p}'\) 。

为证 (b)，只需证明 \(\mathfrak{p}(\mathbf{B})\) 包含在饱和 \(q\) ——理想 \(\mathfrak{pB}\) 关于 \(\mathfrak{p}(\mathbf{B})\) 的——之中（第 II 章，§ 2，第 4 节，命题 10）；由于 \(\mathfrak{p}(\mathbf{B})\) 不包含在任何一个 \(\mathfrak{n}_j\) ( \(1 \leqslant j \leqslant r\) )中，甚至只需证明

\[
\mathfrak {p} (\mathrm{B}) \subset \mathfrak {q} \cup \mathfrak {n} _ {1} \cup \dots \cup \mathfrak {n} _ {r}\tag{4}
\]

即可，这依据第 II 章，§1，第 1 节，命题 2。为此，考虑一个元素 \(u \in \mathfrak{p}(\mathrm{B})\) ，它不属于任何 \(\mathfrak{n}_j\) ( \(1 \leqslant j \leqslant r\) )（第 II 章，§1，第 1 节，命题 2）；设 \(u_1 = u, u_2, \ldots, u_m\) 是 \(u\) 在 \(\mathcal{G}\) 下轨道中的不同元素；记 \(w = u_1u_2\ldots u_m, v = u_2\ldots u_m\) ；显然 \(w \in \mathrm{A}\) ；另一方面，若 \(\tau \in \mathcal{H}\) ，则 \(\tau .u = u\) ，因而必有 \(\tau .u_{i} \neq u\) （对 \(i \geqslant 2\)），这表明 \(\tau .v = v\) ，从而 \(v \in B\) 。与 (a) 的证明一样可以证明：若 \(\sigma \in \mathcal{G}\) 满足 \(\sigma .u = u_{i}\) （其中 \(i \geqslant 2\) ），则 \(\sigma^{-1}.p'\) 位于某个 \(n_{j}\) 之上，且由于 \(u \notin n_{j}\) ，也有 \(u \notin \sigma^{-1}.p'\) ，换言之 \(u_{i} \notin p'\) 。由此断定 \(v \notin p'\) ，因而 \(v \notin p(B)\) 。另一方面，显然 \(w \in p' \cap A = p\) ，而关系式 \(w = uv\) 表明 \(u\) 属于 \(pB\) 关于 \(p(B)\) 的饱和，从而建立了 (4)。

\begin{enumerate}
\def\labelenumi{(\roman{enumi})}
\setcounter{enumi}{2}
\tightlist
\item
  设 \(A'\) 是整环，\(\bigcap_{n\geq0}p^{n}A_{p'}^{\prime}=0\) ，且 (ii) 的条件 (a) 与 (b) 成立。用 (ii) 的同样记号，显然 \(S^{-1}A'\) 是整环且 \(S^{-1}A_{p'}^{\prime}=A_{p'}^{\prime}\) ；因此可以把 \(A'\) 与 \(p'\) 换成 \(S^{-1}A'\) 与 \(S^{-1}p'\) ，换言之，还可设理想 p 是极大的。于是假设 (a) 与 (b) 蕴含
\end{enumerate}

\[
\mathrm{B} _ {\mathfrak {p} (\mathrm{B})} = \mathrm{A} + \mathfrak {p} \mathrm{B} _ {\mathfrak {p} (\mathrm{B})}\tag{5}
\]

对 \(n\) 用归纳法，可导出 \(B_{p(B)} = A + p^n B_{p(B)}\) （对所有 \(n > 0\)）。现设 \(\sigma\) 是 \(\mathcal{G}^{\mathbb{Z}}\) 的元素，\(x\) 是 \(B\) 的元素。对所有 \(n > 0\) ，存在 \(a_n \in A\) 使得 \(x - a_n \in p^n B_{p(B)} \subset p^n A_p'\) ；由于 \(\sigma .a_n = a_n\) 且 \(\sigma .p' = p'\) ，可导出 \(\sigma .x - x \in p^n A_p'\) 。由于这个关系式对所有 \(n\) 成立，由假设便得 \(\sigma .x = x\) 。

\textbf{注}\par

\begin{enumerate}
\def\labelenumi{(\arabic{enumi})}
\setcounter{enumi}{2}
\item
  若 \(\mathbf{A}'\) 是整环且是 Noether 的，则条件 \(\bigcap_{n\geqslant 1}\mathfrak{p}^n\mathrm{A}_{\mathfrak{p}'}' = 0\) 总成立（第 III 章，§3，第 2 节，命题 5 的推论）。可以证明，若设 \(\mathbf{A}'\) 是整环且 \(\mathbf{A}\) 是 Noether 的，该条件也满足。
\item
  若 p 不是 A 的极大理想，则在 (ii) 的假设下关系式 (3) 未必成立，因而即使取 \(H = G^{Z}\) （从而 \(B = A^{Z}\) ），A/p 与 B/p(B) 也未必同构（习题 10）。
\end{enumerate}

推论 1. 在定理 2 的假设下，环 A/p 与 \(\mathbf{A}^{\mathbf{Z}} / (\mathfrak{p}'\cap \mathbf{A}^{\mathbf{Z}})\) 有相同的分式域，且局部环 \((\mathbf{A}^{\mathbf{Z}})_{\mathfrak{p}'\cap \mathbf{A}^{\mathbf{Z}}}\) 的极大理想由 p 生成。

推论 2. 设 A' 是整环，G 是作用于 A' 的有限群，A 是 G 的不变元环，p' 是 A' 的素理想；设 K、\(K^{Z}\) 与 K' 分别是 A、\(A^{Z}\) 与 A' 的分式域。则 K' 是 K 的 Galois 扩张，且 K' 中包含 K 并使 p' 是 A' 的位于 A' ∩ L 的理想 p' ∩ L 之上的唯一素理想的子域 L，恰是包含 \(K^{Z}\) 的那些子域。

\(\mathcal{G}\) 作用于 \(\mathbf{K}'\) ，且 \(\mathbf{K}\) 是 \(\mathcal{G}\) 在 \(\mathbf{K}'\) 中的不变元域（§ 1，第 9 节，命题 23 应用于 \(S = A - \{0\}\) ）；类似地，\(\mathbf{K}^{\mathbf{Z}}\) 是 \(\mathcal{G}^{\mathbf{Z}}\) 的不变元域；于是按定义，\(\mathbf{K}'\) 是 \(\mathbf{K}\) 的 Galois 扩张。若 \(\mathcal{H}\) 是 \(\mathcal{G}\) 中由使 L 的元素都不变的那些 \(\sigma \in \mathcal{G}\) 组成的子群，则说 L 包含 \(K^{\mathbb{Z}}\) 就意味着 \(\mathcal{H}\) 包含于 \(\mathcal{G}^{\mathbb{Z}}\) （《代数》，第 V 章，§ 10，第 5 节，定理 3），且由于 L 是 \(\mathcal{H}\) 在 \(K'\) 中的不变元域，\(A' \cap L\) 是 \(\mathcal{H}\) 在 \(A'\) 中的不变元环；于是第二个断言由命题 4 (i) 得出。

定义 4. 在命题 4 的推论 2 的假设与记号下，若 A' 的位于 p 之上的素理想的个数等于 [K':K]，则称 A 的素理想 p 在 K' 中完全分解。

这等价于说：对 A' 的位于 p 之上的素理想 \(p'\) ，子群 \(\mathcal{G}^{Z}(p')\) 等于使 A' 的所有元素都不变的子群 M；也等价于 \(\mathrm{A}^{\mathbf{Z}}(\mathfrak{p}') = \mathrm{A}'\) ；还等价于 G/M 忠实地作用于 A' 的位于 p 之上的素理想组成的集。

推论 3. 设 A' 是整环，G 是作用于 A' 的有限交换群，A 是 G 的不变元环，p 是 A 的素理想，K 与 K' 分别是 A 与 A' 的分式域。则 A' 的位于 p 之上的素理想都有相同的分解环 \(A^{Z}\) ，且分式域 \(K^{Z}\) （\(A^{Z}\) 的）是 K 与 K' 之间使 p 完全分解的最大中间域。

若 \(\mathfrak{p}'\) 是 A' 的位于 \(\mathfrak{p}\) 之上的素理想，则 \(\mathcal{G}^{\mathbf{Z}}(\sigma. \mathfrak{p}') = \mathcal{G}^{\mathbf{Z}}(\mathfrak{p}')\) ，因为 \(\mathcal{G}\) 交换（公式 (2)），因而（定理 2 (i)）A' 的位于 \(\mathfrak{p}\) 之上的所有素理想有相同的分解群 \(\mathcal{G}^{\mathbf{Z}}\) ，从而有相同的分解环 \(\mathrm{A}^{\mathbf{Z}}\) ；它们的个数是 \((\mathcal{G}: \mathcal{G}^{\mathbf{Z}})\) 。设 L 是 K 与 K' 之间的中间域，\(\mathcal{H}\) 是 \(\mathcal{G}\) 的使 L 的元素不变的子群；\(\mathfrak{p}'\) 关于 \(\mathcal{H}\) 的分解群是 \(\mathcal{G}^{\mathbf{Z}} \cap \mathcal{H}\) ；由于 A' \(\cap\) L 是 \(\mathcal{H}\) 在 A' 中的不变元环，A' 的位于 \(\mathfrak{p}' \cap \mathrm{L}\) 之上的素理想的个数是 \((\mathcal{H}: (\mathcal{G}^{\mathbf{Z}} \cap \mathcal{H})) = (\mathcal{H}\mathcal{G}^{\mathbf{Z}}: \mathcal{G}^{\mathbf{Z}})\) （因为 \(\mathcal{G}\) 交换）。于是 A' \(\cap\) L 的位于 \(\mathfrak{p}\) 之上的素理想的个数是 \((\mathcal{G}: \mathcal{H}\mathcal{G}^{\mathbf{Z}})\) 。因此，为使 \(\mathfrak{p}\) 在 L 中完全分解，必须而且只需 \((\mathcal{G}: \mathcal{H}\mathcal{G}^{\mathbf{Z}}) = [\mathrm{L}: \mathrm{K}] = (\mathcal{G}: \mathcal{H})\) ，而由于 \(\mathcal{H} \subset \mathcal{H}\mathcal{G}^{\mathbf{Z}}\) ，这等价于 \(\mathcal{H}\mathcal{G}^{\mathbf{Z}} = \mathcal{H}\) ，也等价于 \(\mathcal{G}^{\mathbf{Z}} \subset \mathcal{H}\) ，最终等价于 \(L \subset K^{\mathbf{Z}}\) 。

命题 5. 在定理 2 的假设与记号下，分式域 \(k^{\mathrm{T}}\) （\(\mathrm{A}^{\mathrm{T}} / (\mathfrak{p}' \cap \mathrm{A}^{\mathrm{T}})\) 的）等于 \(k_s'\) ，即 \(k\) 的含于 \(k'\) 的最大可分扩张。

与命题 4 一样，这可归结为 p 是 A 的极大理想的情形，此时 \(p'\) 、\(p' \cap A^{Z}\) 与 \(p' \cap A^{T}\) 分别在 \(A'\) 、\(A^{Z}\) 与 \(A^{T}\) 中是极大的（第 1 节，命题 1）。

对所有 \(x \in \mathbf{A}'\) ，多项式 \(\mathbf{P}(\mathbf{X}) = \prod_{\sigma \in \mathcal{G}^{\mathrm{T}}} (\mathbf{X} - \sigma.x)\) 的系数都在惯性环 \(\mathbf{A}^{\mathrm{T}}\) 中，且按 \(\mathcal{G}^{\mathrm{T}}\) 的定义，它在 \(\mathbf{A}'\) 中的全部根都模 \(p'\) 同余；因此多项式 \(\pi'(\mathbf{P})\) ——\(\mathbf{A}^{\mathrm{T}} / (p' \cap \mathbf{A}^{\mathrm{T}})\) 上的，其系数为 \(\mathbf{P}\) 的诸系数在同态 \(\pi': \mathbf{A}' \to \mathbf{A}' / p'\) 下的典范像——在 \(\mathbf{A}' / p'\) 中的全部根都等于 \(x\) 的像，这表明 \(k'\) 是 \(k^{\mathrm{T}}\) 的纯不可分扩张；由此得 \(k_{s}^{\prime} \subset k^{\mathrm{T}}\) ，因为 \(k_{s}^{\prime}\) 的每个元素在 \(k\) 上可分，更不必说在 \(k^{\mathrm{T}}\) 上可分。

我们知道 \(k_s'\) 是 \(k\) 的 Galois 扩张（《代数》，第 V 章，§ 10，第 9 节，命题 14），且由定理 2，其 Galois 群同构于 \(\mathcal{G}' = \mathcal{G}^{\mathrm{Z}} / \mathcal{G}^{\mathrm{T}}\) 。由于 \(k^{\mathrm{T}}\) 是 \(k_s'\) 的纯不可分扩张，\(k^{\mathrm{T}}\) 是 \(k\) 的拟 Galois 扩张，且 \(k^{\mathrm{T}}\) 在 \(k\) 上的次数的可分因子为 \(q = (\mathcal{G}^{\mathrm{Z}}: \mathcal{G}^{\mathrm{T}})\) 。剩下要证 \(k^{\mathrm{T}}\) 是 \(k\) 的可分扩张。上面已经看到，\(\mathcal{G}'\) 等同于 \(A^{\mathrm{T}}\) 的一个自同构群，且 \(A^{\mathrm{Z}}\) 是 \(\mathcal{G}'\) 的不变元环。于是若 \(x \in A^{\mathrm{T}}\) ，则多项式 \(Q(X) = \prod_{\sigma' \in \mathcal{G}'} (X - \sigma'(x))\) 的系数在 \(A^{\mathrm{Z}}\) 中；\(A^{\mathrm{Z}} / (p' \cap A^{\mathrm{Z}})\) 上的、系数为 \(Q\) 的诸系数在 \(\pi'\) 下的像的多项式次数为 \(q\) ，且有一个根 \(\pi'(x) \in A^{\mathrm{T}} / (p' \cap A^{\mathrm{T}})\) 。由于由命题 4 (ii) 有 \(A^{\mathrm{Z}} / (p' \cap A^{\mathrm{Z}}) = k\) ，可见 \(k^{\mathrm{T}}\) 的每个元素的次数都 \(\leqslant q\) （在 \(k\) 上）。

在这样的约定下，设 \(k_{1}\) 是 \(k\) 的拟 Galois 扩张 \(k^{\mathrm{T}}\) 的 \(k\) -自同构群的不变元域；则 \([k^{\mathrm{T}}:k_{1}] = q\) （《代数》，第 V 章，§ 10，第 9 节，命题 14）。设 \(u\) 是 \(k^{\mathrm{T}}\) 在 \(k_{1}\) 上的本原元素；由于它的次数为 \(q\) （在 \(k_{1}\) 上）且次数 \(\leqslant q\) （在 \(k\) 上），故它的次数为 \(q\) （在 \(k\) 上），且它在 \(k_{1}\) 上的极小多项式的系数在 \(k\) 中；这表明 \(u\) 在 \(k\) 上可分。另一方面，对所有 \(v \in k_{1}\) ，存在一个幂 \(p^{f}\) （特征指数 \(p\) 的幂），使得 \(v^{p^{f}} \in k\) 。由此断定 \(k(u - v)\) 包含

\[
(u - v)^{p^{f}} = u^{p^{f}} - v^{p^{f}},
\]

从而包含 \(u^{p^{f}}\) ，因而包含 \(k(u^{p^{f}})\) 。但由于 \(u\) 在 \(k\) 上可分，\(k(u) = k(u^{p^{f}})\) （《代数》，第 V 章，§ 3，第 3 节，命题 4），于是 \(k(u) \subset k(u - v)\) 。由于 \(u\) 的次数为 \(q\) （在 \(k\) 上）而 \(u - v\) 的次数 \(\leqslant q\) ，由此得 \(k(u) = k(u - v)\) ，从而 \(v \in k(u)\) 。这表明 \(v\) 在 \(k\) 上可分，因而 \(k_1 = k\) ，即 \(k^{\mathrm{T}}\) 在 \(k\) 上可分。

推论. 若惯性群 \(\mathcal{G}^{\mathrm{T}}(\mathfrak{p}')\) 的阶与特征指数 \(p\) （\(k\) 的）互素，则域 \(k'\) 是 \(k\) 的 Galois 扩张。

用命题 5 证明中的记号，多项式 \(\pi'(\mathbf{P})\) 的系数在 \(k^{\mathrm{T}} = k_{s}'\) 中，且其全部根都等于 \(\pi'(x)\) ；我们立即导出 \(\pi'(\mathbf{P})\) 是 \(\pi'(x)\) 在 \(k_{s}'\) 上的某个极小多项式的幂；但后者的次数等于 \(p\) 的一个幂，因而由于 \(\pi'(\mathbf{P})\) 的次数等于 \(\mathcal{G}^{\mathrm{T}}\) 的阶，假设蕴含 \(\pi'(x) \in k_{s}'\) ，换言之 \(k_{s}' = k'\) 。
% semantic-fragments: legacy-input-seam
\relax{}
